\section{Evola’s Advice to the Pope}

\begin{quotex}
On the occasion of the visit of Pope Benedict XVI to the USA, it may be appropriate to quote Evola on the situation of the Church just after Vatican II and on its prospects to become an authentic force of Tradition.

This is from the conclusion of the chapter ``Esoteric Catholicism'' in his work \textit{Mask and Face of Contemporary Spirituality}.
\end{quotex}


Nominally, the pontificate, an institution that was already existent in ancient Rome, is part of Catholicism, and stands at the center of the Church and at the top of her hierarchy. But we must ask what subsists in it of its primal function and of Tradition on the whole. The prophetic hope of a Jaochim of Fiore\footnote{\url{https://en.wikipedia.org/wiki/Joachim_of_Fiore}} in the coming of a ``angelic pope'' having almost the marks of an initiate and inaugurating a new ``Reign'', that of the acting, vivifying, and living Holy Spirit, has remained a utopian illusion. And if we want to adapt ourselves to the contingencies of recent times, in particular the figures of the two most recent pontiffs, John XXIII and Paul VI --- with the climate of ``aggiornamento'' and modernization, with the increasing aversion toward Catholic Integralism\footnote{\url{https://en.wikipedia.org/wiki/Integralism}} and toward so-called ``residual medievalism'' --- seem to definitively seal a disastrous balance sheet.

So the conditions for which we can give a somewhat positive response to the issue we formulated, viz., the possibility of Catholicism to furnish what many sought elsewhere, appears to be nonexistent, at least if we consider the large picture. Yet so often this search has pushed them into the confusions and errors of neo-spiritualism. Given what we have said, it is problematic that, in spite of everything, the Church, the ``mystical body of Christ'', is the carrier and administrator of a true supernatural power objectively acting through rites and sacraments, of which it can benefit those who become its members, hoping thereby to experience beyond the confessional religion and not seeing the paramount goal in so-called ``holiness''.

One can nevertheless recognize that Catholicism contains, in spite of everything, traces of a wisdom that can serve as the basis for an ``esoteric'' dimension of various elements. One exponent of ``integral traditionalism'' had this to say:

\begin{quotex}
The fact that the representatives of the Catholic Church understand so little of their own doctrines here he is referring to their inner dimension must not make it such that we ourselves demonstrate the same incomprehension.
\end{quotex}


Otherwise, all the impediments and all the limitations appear difficult to remove. Leaving priestly Catholicism out of consideration, we could refer to the ascetic Catholic in regard, above all, to the ancient monastic traditions, to that which concerns, if not an initiation, at least an interior discipline that looks toward an opening to transcendence and an approach to the supernatural. But even in this case, the difficult work of purification and essentialization would be imposed, given the co-presence of devotional elements and specifically Christian complexes through which it is perhaps difficult to gather up valid tools for interior action without including the knowledge of what other traditions offer.

A Catholicism that is raised up to the level of a truly universal, unanimous, and perennial tradition where faith can be integrated into a metaphysical realization, the symbol integrated into the path to awakening, the rites and sacraments into acts of power, dogmas into expressions of an absolute and infallible consciousness because it is beyond human and, as such, alive in beings unbound from terrestrial chains through an ascesis, and where the pontificate recovers its primal mediating function --- such a Catholicism could supplant every ``spiritualism'', both present and future.

But observing the reality is this anything more than a dream?


\flrightit{Posted on 2008-04-18 by Cologero}

